\paragraph{Context and prior work.}
Haldane's prediction, formulated in 1981 and published in 1983,
distinguishes integer and half-odd-integer isotropic antiferromagnetic
Heisenberg chains, with a positive bulk gap for positive integer spin
\cite{Haldane1981,Haldane1983PLA,Haldane1983PRL}.
For spin one-half, Lieb, Schultz and Mattis had already constructed
excitations of the periodic chain whose energy above the ground state
vanishes with its length \cite[Appendix B]{LSM1961}.
The Affleck--Kennedy--Lieb--Tasaki construction gave a rigorous gapped
spin-one model \cite{AKLT1987,AKLT1988}. Its interaction contains a
biquadratic term: with $q=\mathbf S_j\cdot\mathbf S_{j+1}$, it is
proportional to $q+q^2/3$ up to an additive constant. Our Hamiltonian
has the pure bilinear interaction $q$ and fixed endpoint fields.

The endpoint condition matters even when the bulk interaction is
fixed. White and Huse numerically characterized the effective
spin-one-half degrees of freedom at free ends \cite{WhiteHuse1993}.
To estimate bulk properties, including the gap, they added physical
spin-one-half endpoint sites and adjusted their couplings. Magnetic
fields on the original spin-one endpoints give a different
finite-volume model, whose gap requires its own boundary estimates.

Tasaki formulated precisely this boundary problem in his work on the
topological character of the spin-one Heisenberg ground state
\cite{Tasaki2025}. Equation (2) of that paper is the Hamiltonian
\eqref{eq:target}; Lemma 1 proves uniqueness for every $L\ge1$ and
$h>0$. Assumption 2 asks for a positive gap uniform in all sufficiently
large $L$, at one fixed $h>0$. Theorem~\ref{thm:main} resolves this
fixed-boundary-field Haldane gap assertion positively and establishes
that assumption. Section~\ref{sec:topology} then derives Tasaki index
$-1$ for every boundary-selected subsequential local limit.

The topological viewpoint developed through several distinct
descriptions. Den Nijs and Rommelse introduced nonlocal string order
\cite{DenNijsRommelse1989}, and Kennedy and Tasaki connected it to
hidden symmetry breaking \cite{KennedyTasaki1992}. Pollmann, Turner,
Berg and Oshikawa described symmetry protection through matrix-product
states and the entanglement spectrum \cite{Pollmann2010}; Ogata
extended the bond-centered reflection index to infinite-chain states
beyond the matrix-product description \cite{Ogata2021}.
Tasaki's local-twist index \cite{Tasaki2018} instead uses the twist
expectation studied by Nakamura and Todo \cite{NakamuraTodo2002}.
The boundary-state theorem in \cite{Tasaki2025} uses uniform $z$
rotations and site-centered reflection. These invariances imply
invariance under site reflection combined with a global $\pi$ rotation
about $z$, the symmetry studied by Fuji, Pollmann and Oshikawa
\cite{FujiPollmannOshikawa2015}.
Our consequence concerns Tasaki's index in this precise symmetry
setting; it does not identify the different diagnostics above.

\paragraph{Methodological inputs.}
For frustration-free chains, Lemm and Mozgunov proved finite-size
criteria involving both interior and boundary-segment gaps
\cite{LemmMozgunov2019}. Here the boundary information comes from
weighted spatial moments and a coupled Gibbs-purity recurrence.
The finite trial vectors use the matrix-product framework of finitely
correlated states \cite{FNW1992} and its finite-chain formulation
\cite{PerezGarcia2007}; they certify energy upper bounds that
initialize the recurrence.

The continuous-history representation is an instance of the
time-ordered Poisson expansion of Aizenman and Nachtergaele
\cite{AizenmanNachtergaele1994}. The particular spatial transfer
operator and finite periodic estimates used here come from
\cite{periodic}, which proves a uniform gap on even periodic rings.
We retain that source's spin normalization, coupling, and energy
shift, and state explicitly the transfer and finite-input results
used below. The endpoint pairings and all estimates specific to the
open chain are proved here.
